% ============================================================
%  Setmana 3 — Unitat 7 · Límits i continuïtat
%  Exercicis: 92, 93, 94, 102, 103, 40, 106
% ============================================================

\setmana{3}{Unitat 7 \textperiodcentered{} Límits i continuïtat}{4 d'octubre de 2026}%
        {92, 93, 94, 102, 103, 40 i 106}

\begin{errada}{Exercici 92 b)}
  \ondiu{$\dlim_{x\to 0^-} f(x)=\dlim_{x\to 0^+} f(x)=0{,}5$ \quad i \quad $f(2)=2{,}5$}
  \hadedir{$\dlim_{x\to 0^-} f(x)=\dlim_{x\to 0^+} f(x)=1$ \quad i \quad $f(2)=5$}
  \motiu{A la gràfica de l'enunciat cada quadre és una unitat (a l'eix $Y$,
    el $2$ és a la segona marca): el punt buit és a alçada $1$ i la corba
    passa pel punt $(2,5)$. És la mateixa paràbola de l'apartat d), on el
    solucionari sí que llegeix $f(-1)=2$ i límit $5$ en $x=2$.
    La conclusió no canvia: discontinuïtat evitable en $x=0$ i funció
    contínua en $x=2$.}
\end{errada}

\begin{errada}{Exercici 92 e)}
  \ondiu{«La funció no és contínua en $x=-2$; té una discontinuïtat de salt finit.»}
  \hadedir{«La funció no és contínua en $x=2$; té una discontinuïtat de salt finit.»}
  \motiu{És la frase que tanca l'estudi del punt $x=2$, on $f(2)=3$ i els
    límits laterals valen $-1$ i $3$. El punt $x=-2$ ja s'ha estudiat
    abans, i hi ha un salt infinit.}
\end{errada}

\begin{errada}{Exercici 93 c)}
  \ondiu{contínua en el seu domini, $(-\infty,-2)\cup(2,+\infty)$}
  \hadedir{contínua en el seu domini, $(-\infty,-2]\cup[2,+\infty)$}
  \motiu{La condició és $x^2-4\geq 0$, que també es compleix en $x=-2$ i en
    $x=2$: allà la funció existeix i val $\sqrt{0}=0$. La mateixa solució
    ho escriu bé just abans ($x\geq 2$, $x\leq -2$).}
\end{errada}

\begin{errada}{Exercici 93 g)}
  \ondiu{$\dlim_{x\to 1} 2\,|x-1| = 0 = y(0)$}
  \hadedir{$\dlim_{x\to 1} 2\,|x-1| = 0 = y(1)$}
  \motiu{El punt que s'estudia és $x=1$, i $y(1)=2\,|1-1|=0$.
    En canvi, $y(0)=2$.}
\end{errada}

\begin{errada}{Exercici 94 b)}
  \ondiu{$x^2-2x=3\neq 0$ per a qualsevol $x$ real}
  \hadedir{$x^2-2x+3\neq 0$ per a qualsevol $x$ real}
  \motiu{El denominador és $x^2-2x+3$, i no s'anul·la mai perquè el
    discriminant és negatiu: $(-2)^2-4\cdot 1\cdot 3=-8<0$.
    Per això la funció és contínua en tot $\R$.}
\end{errada}

\begin{errada}{Exercici 94 e), punt $x=-3$}
  \ondiu{«No existeix $\dlim_{x\to 3} f(x)$.»}
  \hadedir{«No existeix $\dlim_{x\to -3} f(x)$.»}
  \motiu{Els dos límits laterals que s'acaben de calcular són els del punt $x=-3$.}
\end{errada}

\begin{errada}{Exercici 94 e), punt $x=1$}
  \ondiu{«La funció té una discontinuïtat evitable en $x=3$.»}
  \hadedir{«La funció té una discontinuïtat evitable en $x=1$.»}
  \motiu{El límit que s'ha calculat és el de $x\to 1$ (val $\tfrac{1}{8}$), i
    $x=1$ és el zero del denominador que se simplifica.
    En $x=3$ la funció és contínua.}
\end{errada}

\begin{errada}{Exercici 102}
  \ondiu{$\dlim_{x\to -1^-} f(x)=5$}
  \hadedir{$\dlim_{x\to -1^-} f(x)=-5$}
  \motiu{Per a $x<-1$ la funció és $3x-2$, i $3\cdot(-1)-2=-5$.
    La conclusió no canvia: els límits laterals ($-5$ i $-4$) són
    diferents i en $x=-1$ hi ha una discontinuïtat de salt finit.}
\end{errada}

%\senseerrades{103, 40 i 106}
